% Frozen counterexample candidate, 2026-08-31.
% Requires independent mathematical audit and production verification.
% Source abbreviations and precise versions are specified in SOURCE_INTERFACE.md.
\begin{proof}
We use the model-metric convention of Gubler--Martin [GM19,
Introduction and Sections 2--4]: a \(\mathbf Q\)-formal metric becomes a
formal metric after a positive integral tensor power, and semipositivity
means that a representing formal line bundle is nef on the special fibre.
All slopes below are measured using the intrinsic logarithmic annulus
length, not an arbitrarily rescaled coordinate on an edge.

\emph{Step 1: the field, the Tate curve, and its marked skeleton.}
For the field example in the statement, the weighted Gauss norm on
\(\mathbf Q_p[T]\) is
\[
 \left|\sum_i c_iT^i\right|
   =\max_i |c_i|_p p^{-\alpha i}.
\]
Distinct nonzero monomials have distinct valuations because
\(\alpha\notin\mathbf Q\); hence the term of largest norm is unique,
and the same observation proves multiplicativity. Extend the norm to
\(\mathbf Q_p(T)\), complete, and take a completed algebraic closure.
The resulting complete algebraically closed non-Archimedean field
contains \(p,T\) with valuations \(\log p,\alpha\log p\).

Tate's construction applies to every complete nontrivially valued field
and every \(q\) with \(0<|q|<1\): the Weierstrass equation and
nonvanishing discriminant in [Tate95, equations (13)--(18)] give the
elliptic curve, and [Tate95, Theorem 1] gives the multiplicative
uniformization on rational points. Its analytic uniformization and
skeletal description are [BPR14, Remark 4.24]. In particular, the minimal
skeleton is
\[
 \Sigma=\mathbf R/\ell\mathbf Z,
\]
with coordinate induced by \(v(z)\). Its retraction \(\tau\) sends
\(O\) to \(x_0=0\) and \(P\) to \(x_a=a\).
The two points \(x_0,x_a\) have type 2, since their coordinates belong
to \(v(K^\times)\). By [BPR14, Proposition 3.13(2) and Theorem 4.11],
\(\{x_0,x_a\}\) is a semistable vertex set corresponding to a proper
semistable formal model \(\mathfrak E\). Its two edges have lengths
\(a\) and \(\ell-a\); neither is a loop edge, so the model is strictly
semistable. There are two special-fibre components \(C_0,C_a\), and
all component multiplicities and nodal residue-field degrees are one.
If algebraic models are preferred, the formal model algebraizes as in
[GJKM19, Remark 3.3].

The points \(O\) and \(P\) reduce to smooth points of \(C_0\) and
\(C_a\), respectively. Indeed, a node has open-annular tube whose
skeleton is the interior of an edge, whereas a residue disc at a smooth
point retracts to the vertex; this is the formal-fibre description in
[BPR14, Section 1.3 and Proposition 4.10]. Concretely, the reductions
are the unit point in the annular coordinates \(z\) and \(z/u\), not
either of the nodal directions \(0,\infty\).
Their horizontal closures are Cartier divisors: in a smooth formal
coordinate at either section the closure is cut out by that coordinate
minus the coordinate of the section. Reduction gives a simple
rational point on the indicated component and meets no other component.
Thus
\[
 \mathfrak L_0
   :=\mathcal O_{\mathfrak E}(\overline P-\overline O)
\]
is a formal model of \(L\) with component degrees
\begin{equation}
\label{eq:tate-reference-degrees}
 \deg(\mathfrak L_0|C_0)=-1,
 \qquad \deg(\mathfrak L_0|C_a)=1.
\end{equation}
Equivalently its model curvature measure is
\(\delta_{x_a}-\delta_{x_0}\).
This last formulation uses only the usual curve model-curvature
formula; the rest of the obstruction can instead be read entirely in
terms of the two component degrees in
\eqref{eq:tate-reference-degrees}.
Since \(P,O\) are \(K\)-points, \(\deg L=1-1=0\).

\emph{Step 2: an arbitrary hypothetical nef model can be tested on a
common semistable refinement.}
Suppose that some \(m>0\) and some proper admissible formal model
carry a nef model \(\mathfrak N\) of \(L^{\otimes m}\).
Let \(\|\cdot\|_N\) be its formal metric and
\(\|\cdot\|_0\) the metric of \(\mathfrak L_0\), and put
\[
 \varphi=-\log\frac{\|\cdot\|_N}{\|\cdot\|_0^{\otimes m}}.
\]
This is a \(\mathbf Z\)-model function; the metric--function dictionary
is [GM19, Section 4.7]. Take a common model and then a strictly
semistable model \(\mathfrak E'\) dominating \(\mathfrak E\) on which
\(\varphi\) is determined. The semistable domination used here is the
one recalled in the proof of [GJKM19, Proposition 3.7], citing
Bosch--L\"utkebohmert [BL85, Section 7]; it applies over our already
algebraically closed complete field. Formal/algebraic interchange on a
proper curve is allowed by [GM19, Theorem 1.1] and [GJKM19, Remark 3.3].
Nefness is preserved by pullback to the common model, by the projection
formula; alternatively [GM19, Proposition 3.5] tests semipositivity on
any representing model.

Write \(\Sigma'=S(\mathfrak E')\) and let
\(\tau':E_q^{\mathrm{an}}\to\Sigma'\) be its retraction.
The skeleton \(\Sigma'\) is obtained by subdividing \(\Sigma\) and
attaching finitely many finite trees, each at one point of \(\Sigma\).
This follows from [BPR14, Proposition 3.13] and the open-ball components
of the complement of the minimal skeleton. By
[GJKM19, Proposition 3.5],
\[
 \varphi=F\circ\tau',
\]
where \(F\) is continuous, affine on the edges of a finite subdivision
of \(\Sigma'\), and all its slopes are integers. Additional subdivision
points can be incorporated in the semistable vertex set.
Pullback of \(\mathfrak L_0^{\otimes m}\) retains degree \(-m\) and
\(m\) at the vertices \(x_0,x_a\); its degree on every new component
is zero. In fact a new component maps to a closed point of the original
special fibre, and the old components map birationally to the old ones.

\emph{Step 3: positivity leaves no curvature on a degree-zero graph.}
Use the outgoing-slope convention of [GJKM19, Proposition 3.7]. For
each vertex \(x\) of the refined graph set
\[
 b_x=m\,\mathbf 1_{\{x_a\}}(x)-m\,\mathbf 1_{\{x_0\}}(x).
\]
The slope formula and nefness of \(\mathfrak N\) give
\begin{equation}
\label{eq:tate-nef-vertex-inequality}
 b_x+\sum_{\nu\in T_x\Sigma'}\partial_\nu F\ \geq\ 0.
\end{equation}
All tangent-direction weights are one because the residue field is
algebraically closed. Summing the left-hand sides over all vertices
gives zero: \(\sum_x b_x=0\), and each edge contributes its slope and
its negative at its two endpoints. A finite sum of nonnegative numbers
is zero only if every summand is zero. Therefore
\begin{equation}
\label{eq:tate-zero-vertex-curvature}
 b_x+\sum_{\nu\in T_x\Sigma'}\partial_\nu F=0
 \quad\hbox{at every vertex }x.
\end{equation}
No global CLD-current positivity theorem or approximation theorem is
needed for this finite-graph conclusion.

On every tree attached to the circle, \(F\) is constant. To prove this,
take a leaf away from the attachment point. Here \(b_x=0\), and its
single outgoing slope is zero by
\eqref{eq:tate-zero-vertex-curvature}. Remove that edge and leaf and
repeat. This terminates because the tree is finite. In particular,
attached trees contribute no slope at their roots.
At a subdivision vertex of either original circle edge,
\eqref{eq:tate-zero-vertex-curvature} now says that the two oriented
slopes agree. Hence the restriction of \(F\) to the original circle
has a constant integer slope \(s\) on \((0,a)\) and a constant integer
slope \(t\) on \((a,\ell)\), with the same orientation of the
coordinate on both intervals.

At \(x_0\) the two outgoing slopes are \(s\) and \(-t\), and its
reference degree is \(-m\). Thus
\begin{equation}
\label{eq:tate-jump}
 s-t=m.
\end{equation}
The continuity and single-valuedness of \(F\) on the circle give
\begin{equation}
\label{eq:tate-period}
 as+(\ell-a)t=0.
\end{equation}
Solving \eqref{eq:tate-jump}--\eqref{eq:tate-period} yields
\[
 t=-m\frac{a}{\ell},\qquad
 s=m\left(1-\frac{a}{\ell}\right).
\]
Both numbers are irrational, contradicting their integrality. This
rules out the arbitrary model \(\mathfrak N\), for every \(m>0\).
Equivalently, working directly with a \(\mathbf Q\)-model correction
would require the rational slopes
\(t=-a/\ell\) and \(s=1-a/\ell\), which is the same contradiction.
Since a numerically trivial model is nef, the asserted nonexistence of
numerically trivial \(\mathbf Q\)-formal metrics follows as well.

\emph{Step 4: the canonical metric satisfies the CLD hypothesis.}
An elliptic curve is an algebraic abeloid variety, and our ground-field
absolute value is nontrivial. Fix any rigidification of \(L\) at the
origin. [CLD25, Lemma 19.1.6] supplies its continuous canonical cubist
metric. On an algebraic curve numerical triviality of a line bundle is
equivalent to degree zero, so both \(L\) and \(L^{-1}\) are nef.
[CLD25, Proposition 19.3.1] states, under a nontrivial absolute value,
that for an algebraic abeloid the cubist metric is psh if and only if
the underlying line bundle is nef. [CLD25, Corollary 19.3.2] states
that its curvature current vanishes if and only if the line bundle is
numerically trivial. Applied here, these give
\[
 \|\cdot\|_{\mathrm{cub}}\ \text{locally CLD-psh},
 \qquad c_1(L,\|\cdot\|_{\mathrm{cub}})=0.
\]
The split-uniformization setting used in those results is satisfied
without any field extension: for a Tate curve the torus is
\(\mathbf G_m\), the good-reduction abeloid is a point, and the lattice
is \(q^{\mathbf Z}\).

For completeness the stronger local regularity asserted in the
statement follows from the same source. By [CLD25, Proposition 19.2.5],
after trivializing the pullback of \(L\) to \(\mathbf G_m^{\mathrm{an}}\),
the cubist weight differs from a constant by a quadratic function of
the tropical coordinate, with quadratic part \(q_L/2\).
[CLD25, Corollary 19.3.2] gives \(q_L=0\) for our numerically trivial
line bundle, so this weight is real affine in \(-\log|z|\).
The latter is a CLD-smooth function furnished by an analytic unit and
has zero \(d'd''\). The Tate covering has local analytic sections, so
these weights descend locally. Changing an invertible local frame
adds the logarithm of the absolute value of an analytic unit, again
smooth and harmonic. Thus the metric satisfies the local CLD-psh
hypothesis in the ordinary analytic topology, not just a formal
statement about its global curvature mass.

\emph{Step 5: the fixed-line-bundle closure is empty.}
In the convention of [GM19, Introduction], a continuous semipositive
metric on \(L\) is a uniform limit of semipositive model metrics on
the same \(L\). Step 3 shows that there are no such model metrics at
all. Hence there cannot be an approximating sequence for the metric
constructed in Step 4. This proves the stated obstruction, without
making any assertion about definitions that allow an ample twist,
varying generic-fibre line bundles, or larger cones of real local
potentials.
\end{proof}
